% This file aims at providing detailed explanations, detailed and step-by-step mathematical derivations of the concepts in `cgenarris_layer_generation_algorithm_update.md`.
% All notations and symbols should be clearly explained.

\documentclass[11pt]{article}

\usepackage[margin=1in]{geometry}
\usepackage{amsmath}
\usepackage{amssymb}
\usepackage{bm}
\usepackage{enumitem}

\title{Supplementary Notes for the cgenarris Layer-Generation Algorithm}
\author{}
\date{}

\newcommand{\R}{\mathbb{R}}
\newcommand{\dd}{\mathrm{d}}
\newcommand{\T}{\mathsf{T}}

\begin{document}

\maketitle

\section{Scope of This Supplement}

This supplement records detailed explanations for the molecular point-group
detection part of the new layer-generation algorithm.  It currently covers only
the concepts discussed so far:
\begin{enumerate}[label=(\roman*)]
    \item atom equivalence classes;
    \item the gyration tensor and its principal axes;
    \item why symmetry axes are related to eigenvectors of the gyration tensor;
    \item candidate rotation axes and mirror-plane normals;
    \item proper rotations, mirror reflections, and atom permutations;
    \item operation identity by \((\pi,\det M)\), exact group closure, and
    permutation composition;
    \item nonlinear versus linear molecules;
    \item orthogonal Procrustes matrix refinement;
    \item recording symmetry elements and classifying the point group.
\end{enumerate}

\section{Notation and Coordinate Convention}

Consider a molecule with \(N\) atoms.  Before detecting molecular symmetry, the
coordinates are shifted so that the centroid is at the origin.  If
\(\mathbf r_i^{\,0}\in\R^3\) is the input Cartesian coordinate of atom \(i\),
then the molecular centroid is
\[
    \bar{\mathbf r}
    =
    \frac{1}{N}\sum_{i=1}^{N}\mathbf r_i^{\,0},
\]
and the centered coordinate of atom \(i\) is
\[
    \mathbf r_i
    =
    \mathbf r_i^{\,0} - \bar{\mathbf r}.
\]
Here:
\begin{itemize}
    \item \(i\in\{1,\dots,N\}\) is an atom index;
    \item \(\mathbf r_i\in\R^3\) is a column vector
    \[
        \mathbf r_i =
        \begin{bmatrix}
        x_i\\ y_i\\ z_i
        \end{bmatrix};
    \]
    \item \(x_i,y_i,z_i\) are Cartesian coordinates measured relative to the
    molecular centroid;
    \item the superscript \({\T}\) denotes transpose, so
    \(\mathbf r_i^{\T}\) is a row vector;
    \item \(\tau\) denotes the user-specified distance tolerance for symmetry
    matching.
\end{itemize}

All point-group operations in this part of the algorithm act on these centered
Cartesian coordinates.  No lattice vectors or fractional coordinates are used in
molecular point-group detection.

\section{Atom Equivalence Classes}

\subsection{Definition}

Each atom has a chemical element label, denoted \(e_i\), and a centered Cartesian
coordinate \(\mathbf r_i\).  For atoms \(i\) and \(k\), define the interatomic
distance
\[
    d_{ik} = \lVert \mathbf r_i - \mathbf r_k\rVert .
\]
The distance-environment signature of atom \(i\) is the sorted multiset
\[
    \mathrm{sig}(i)
    =
    \big\{(e_k,d_{ik}) : k\ne i\big\}.
\]
That is, \(\mathrm{sig}(i)\) records the element and distance of every other atom
as seen from atom \(i\).  The entries are sorted so that two signatures can be
compared entry by entry.

Two atoms \(i\) and \(j\) are placed in the same equivalence class if:
\begin{enumerate}[label=(\alph*)]
    \item they have the same element, \(e_i=e_j\); and
    \item their sorted distance-environment signatures match within the
    tolerance \(\tau\).
\end{enumerate}
In words, two equivalent atoms are indistinguishable by their element type and by
their full distance environment in the molecule.

\subsection{Why Equivalence Classes Are Useful}

Let \(Q\) be a true molecular symmetry operation.  In this section, \(Q\) is a
\(3\times 3\) orthogonal Cartesian matrix, such as a rotation, reflection,
inversion, or improper rotation.  If \(Q\) is a symmetry of the molecule, then it
maps each atom onto another atom of the same element.  Therefore there exists a
permutation \(\pi\) of the atom indices such that
\[
    Q\mathbf r_i = \mathbf r_{\pi(i)}
\]
for every atom \(i\), up to the tolerance \(\tau\).  The notation \(\pi(i)\)
means ``the index of the atom onto which atom \(i\) is mapped.''

Because \(Q\) is orthogonal, it preserves distances:
\[
    \lVert Q\mathbf r_i - Q\mathbf r_k\rVert
    =
    \lVert \mathbf r_i - \mathbf r_k\rVert .
\]
Using \(Q\mathbf r_i=\mathbf r_{\pi(i)}\) and
\(Q\mathbf r_k=\mathbf r_{\pi(k)}\), this becomes
\[
    \lVert \mathbf r_{\pi(i)} - \mathbf r_{\pi(k)}\rVert
    =
    \lVert \mathbf r_i - \mathbf r_k\rVert .
\]
Thus a true symmetry operation preserves both element labels and all
interatomic distances.  Consequently, if two atoms are related by a true
molecular symmetry operation, they must lie in the same equivalence class.

The algorithm uses this fact in two ways:
\begin{enumerate}[label=(\alph*)]
    \item \textbf{Restricted atom matching.}  When verifying a trial symmetry
    operation, atom \(i\) only needs to be matched against atoms in its own
    equivalence class.
    \item \textbf{Candidate direction generation.}  Candidate symmetry axes and
    mirror normals are constructed from vectors involving atoms in the same
    equivalence class.  This focuses the search on geometrically plausible
    symmetry directions.
\end{enumerate}

\subsection{Completeness of the Equivalence-Class Test}

The equivalence-class criterion is complete as a necessary condition: if two
atoms are related by an exact symmetry operation, then they have the same
element and the same distance environment.  With a sufficiently large tolerance
\(\tau\), the same remains true for slightly relaxed numerical geometries.

The criterion is not sufficient: two atoms may accidentally have identical
distance environments even if no actual symmetry operation maps one atom to the
other.  This does not create an incorrect symmetry operation, because every
candidate operation is later verified explicitly by applying its matrix to all
atoms and checking whether a one-to-one atom mapping exists within the tolerance
\(\tau\).

\section{The Gyration Tensor and Its Principal Axes}

\subsection{Definition of the Gyration Tensor}

The unweighted gyration tensor used here is
\[
    T = \sum_{i=1}^{N} \mathbf r_i \mathbf r_i^{\T}.
\]
Here:
\begin{itemize}
    \item \(T\) is a \(3\times 3\) real matrix;
    \item \(\mathbf r_i\) is the centered Cartesian coordinate of atom \(i\);
    \item \(\mathbf r_i^{\T}\) is the transpose of \(\mathbf r_i\);
    \item \(\mathbf r_i\mathbf r_i^{\T}\) is the outer product of
    \(\mathbf r_i\) with itself;
    \item the sum adds one such \(3\times 3\) matrix for every atom.
\end{itemize}
For one atom,
\[
    \mathbf r_i\mathbf r_i^{\T}
    =
    \begin{bmatrix}
    x_i x_i & x_i y_i & x_i z_i\\
    y_i x_i & y_i y_i & y_i z_i\\
    z_i x_i & z_i y_i & z_i z_i
    \end{bmatrix}.
\]
Therefore the entries of \(T\) are
\[
    T_{\alpha\beta}
    =
    \sum_{i=1}^{N}(\mathbf r_i)_\alpha(\mathbf r_i)_\beta,
\]
where \(\alpha,\beta\in\{x,y,z\}\), and \((\mathbf r_i)_\alpha\) denotes the
\(\alpha\)-component of atom \(i\)'s centered coordinate.

\subsection{Why \(T\) Is Symmetric}

The transpose of one outer product is
\[
    \left(\mathbf r_i\mathbf r_i^{\T}\right)^{\T}
    =
    \mathbf r_i\mathbf r_i^{\T}.
\]
Thus every term in the sum is symmetric.  Equivalently,
\[
    T_{\alpha\beta}
    =
    \sum_i(\mathbf r_i)_\alpha(\mathbf r_i)_\beta
    =
    \sum_i(\mathbf r_i)_\beta(\mathbf r_i)_\alpha
    =
    T_{\beta\alpha},
\]
because scalar multiplication commutes.  Hence \(T=T^{\T}\).

\subsection{Physical Meaning: Squared Molecular Spread}

Let \(\mathbf u\in\R^3\) be a unit vector, \(\lVert\mathbf u\rVert=1\).
The scalar projection of atom \(i\) along direction \(\mathbf u\) is
\[
    \mathbf u\cdot \mathbf r_i = \mathbf u^{\T}\mathbf r_i.
\]
The total squared spread of the molecule along direction \(\mathbf u\) is
\[
    q(\mathbf u)
    =
    \sum_{i=1}^{N}(\mathbf u^{\T}\mathbf r_i)^2.
\]
This can be written in terms of \(T\):
\[
\begin{aligned}
    q(\mathbf u)
    &=
    \sum_{i=1}^{N}
    (\mathbf u^{\T}\mathbf r_i)(\mathbf r_i^{\T}\mathbf u)\\
    &=
    \mathbf u^{\T}
    \left(\sum_{i=1}^{N}\mathbf r_i\mathbf r_i^{\T}\right)
    \mathbf u\\
    &=
    \mathbf u^{\T}T\mathbf u.
\end{aligned}
\]
Thus \(\mathbf u^{\T}T\mathbf u\) has a direct geometric meaning: it measures
the molecule's total squared extent along the direction \(\mathbf u\).

\subsection{Why Principal Axes Are Eigenvectors of \(T\)}

The natural geometric directions are the unit vectors \(\mathbf u\) for which
the squared spread \(q(\mathbf u)=\mathbf u^{\T}T\mathbf u\) is stationary.
These are obtained by optimizing \(q(\mathbf u)\) subject to the unit-length
constraint
\[
    \mathbf u^{\T}\mathbf u = 1.
\]
Introduce a Lagrange multiplier \(\lambda\) and define
\[
    \mathcal L(\mathbf u,\lambda)
    =
    \mathbf u^{\T}T\mathbf u
    -
    \lambda(\mathbf u^{\T}\mathbf u-1).
\]
Taking the derivative with respect to \(\mathbf u\) gives
\[
    2T\mathbf u - 2\lambda \mathbf u = 0.
\]
Therefore
\[
    T\mathbf u = \lambda \mathbf u.
\]
This is precisely the eigenvector equation for \(T\).  Hence the principal
directions of molecular spread are eigenvectors of the gyration tensor.

The eigenvalue \(\lambda\) is the squared spread along the corresponding
eigenvector.  The largest eigenvalue gives the direction of largest second
moment, the smallest eigenvalue gives the direction of smallest second moment,
and the remaining eigenvalue gives the intermediate direction.  This is why the
principal axes can be interpreted as long, medium, and short axes in the
second-moment sense.

\subsection{Orthogonality and Degeneracy}

Because \(T\) is a real symmetric matrix, the spectral theorem guarantees that
\(T\) has three real eigenvalues and admits an orthonormal eigenbasis.  Thus one
can always choose three mutually orthogonal eigenvectors of \(T\).

If two eigenvalues are distinct, their eigenvectors are automatically
orthogonal.  To see this, suppose
\[
    T\mathbf u = \lambda \mathbf u,
    \qquad
    T\mathbf v = \mu \mathbf v,
    \qquad
    \lambda\ne\mu.
\]
Then
\[
    \lambda \mathbf u^{\T}\mathbf v
    =
    (T\mathbf u)^{\T}\mathbf v
    =
    \mathbf u^{\T}T^{\T}\mathbf v
    =
    \mathbf u^{\T}T\mathbf v
    =
    \mu \mathbf u^{\T}\mathbf v.
\]
Therefore
\[
    (\lambda-\mu)\mathbf u^{\T}\mathbf v=0.
\]
Since \(\lambda\ne\mu\), it follows that
\[
    \mathbf u^{\T}\mathbf v=0.
\]
Thus \(\mathbf u\) and \(\mathbf v\) are orthogonal.

If an eigenvalue is degenerate, the corresponding eigenspace has dimension
larger than one.  In that case the individual eigenvectors inside the degenerate
subspace are not unique, although an orthonormal basis can always be chosen
inside that subspace.  For example:
\begin{itemize}
    \item if \(\lambda_1>\lambda_2>\lambda_3\), all three principal axes are
    unique up to sign;
    \item if \(\lambda_1=\lambda_2\ne\lambda_3\), the nondegenerate axis is
    unique, but any orthonormal pair spanning the two-dimensional degenerate
    plane is equally valid;
    \item if \(\lambda_1=\lambda_2=\lambda_3\), the tensor is isotropic and does
    not identify any preferred direction.
\end{itemize}
This is why the algorithm does not rely only on the three eigenvectors of
\(T\).  It also adds atom-derived candidate directions, which are essential in
degenerate or nearly degenerate cases.

\subsection{Why a Nondegenerate Symmetry Axis Must Be a Principal Axis}

Let \(Q\) be a true molecular symmetry operation represented by a \(3\times 3\)
orthogonal matrix.  Since \(Q\) is a symmetry, it permutes the atoms.  Therefore
there exists a permutation \(\pi\) such that
\[
    Q\mathbf r_i = \mathbf r_{\pi(i)}
\]
for all atoms \(i\).

Now transform the gyration tensor:
\[
\begin{aligned}
    QTQ^{\T}
    &=
    Q
    \left(\sum_{i=1}^{N}\mathbf r_i\mathbf r_i^{\T}\right)
    Q^{\T}\\
    &=
    \sum_{i=1}^{N}
    (Q\mathbf r_i)(Q\mathbf r_i)^{\T}\\
    &=
    \sum_{i=1}^{N}
    \mathbf r_{\pi(i)}\mathbf r_{\pi(i)}^{\T}.
\end{aligned}
\]
Because \(\pi\) is a permutation, the final sum is just a reordering of the
original atom sum.  Hence
\[
    QTQ^{\T} = T.
\]
Multiplying this equation on the right by \(Q\), and using \(Q^{\T}Q=I\), gives
\[
    QT = TQ.
\]
Thus \(T\) commutes with every true molecular symmetry operation \(Q\).

Now suppose \(Q\) is a nontrivial proper rotation about the unit axis
\(\mathbf a\).  The axis direction is fixed by the rotation:
\[
    Q\mathbf a = \mathbf a.
\]
Using \(QT=TQ\),
\[
    Q(T\mathbf a)
    =
    T(Q\mathbf a)
    =
    T\mathbf a.
\]
Therefore \(T\mathbf a\) is also fixed by the rotation \(Q\).  For a nontrivial
proper rotation in three dimensions, the only fixed vectors are those parallel
to the rotation axis \(\mathbf a\).  Hence \(T\mathbf a\) must be parallel to
\(\mathbf a\), so there exists a scalar \(\lambda\) such that
\[
    T\mathbf a = \lambda\mathbf a.
\]
This is the eigenvector equation.  Therefore the rotation axis
\(\mathbf a\) is an eigenvector of \(T\).

If the eigenvalues of \(T\) are nondegenerate, the eigenvector directions are
unique up to sign.  Consequently, any true proper rotation axis must coincide
with one of the three principal axes of \(T\).  If the eigenvalues are
degenerate, the tensor may identify only a subspace, not a unique axis, and
additional atom-derived candidate directions are required.

\section{Candidate Symmetry Directions}

\subsection{Sources of Candidate Directions}

The algorithm builds a finite list of candidate directions.  A direction may
later be interpreted either as a rotation axis or as a mirror-plane normal,
depending on the trial operation being tested.

The candidate set contains:
\begin{enumerate}[label=(\alph*)]
    \item the three principal axes of the gyration tensor \(T\);
    \item for atoms \(i,j\) in the same equivalence class, the directions
    \[
        \mathbf r_i,\qquad
        \mathbf r_i+\mathbf r_j,\qquad
        \mathbf r_i-\mathbf r_j,\qquad
        \mathbf r_i\times\mathbf r_j.
    \]
\end{enumerate}
Zero or near-zero vectors are discarded.  All remaining directions are
normalized to unit length and deduplicated up to sign, since the geometric line
defined by \(\mathbf a\) is the same as that defined by \(-\mathbf a\).

\subsection{Geometric Motivation}

These atom-derived directions cover common symmetry geometries:
\begin{itemize}
    \item \(\mathbf r_i\) captures axes or normals passing through an atom and
    the centroid;
    \item \(\mathbf r_i+\mathbf r_j\) captures bisectors of equivalent atom
    pairs;
    \item \(\mathbf r_i-\mathbf r_j\) captures the normal direction for a mirror
    plane that swaps atoms \(i\) and \(j\);
    \item \(\mathbf r_i\times\mathbf r_j\) captures a direction normal to the
    plane spanned by \(\mathbf r_i\) and \(\mathbf r_j\).
\end{itemize}
The principal axes supply global second-moment directions.  The atom-derived
directions supply additional candidates in degenerate cases where the gyration
tensor does not determine all symmetry axes uniquely.

\section{Trial Symmetry Operation Matrices}

\subsection{Proper Rotations}

Let \(\mathbf a\in\R^3\) be a unit candidate rotation axis, and let \(\theta\)
be a trial rotation angle.  The rotation matrix about \(\mathbf a\) by angle
\(\theta\) is given by Rodrigues' formula:
\[
    R(\mathbf a,\theta)
    =
    \cos\theta\, I
    +
    \sin\theta\, [\mathbf a]_\times
    +
    (1-\cos\theta)\,\mathbf a\mathbf a^{\T}.
\]
Here:
\begin{itemize}
    \item \(R(\mathbf a,\theta)\) is a \(3\times 3\) rotation matrix;
    \item \(I\) is the \(3\times 3\) identity matrix;
    \item \(\mathbf a\mathbf a^{\T}\) is the projection matrix onto the axis
    direction \(\mathbf a\);
    \item \([\mathbf a]_\times\) is the skew-symmetric cross-product matrix
    satisfying
    \[
        [\mathbf a]_\times \mathbf v = \mathbf a\times\mathbf v
    \]
    for any vector \(\mathbf v\).
\end{itemize}
If
\[
    \mathbf a =
    \begin{bmatrix}
    a_x\\ a_y\\ a_z
    \end{bmatrix},
\]
then
\[
    [\mathbf a]_\times
    =
    \begin{bmatrix}
    0 & -a_z & a_y\\
    a_z & 0 & -a_x\\
    -a_y & a_x & 0
    \end{bmatrix}.
\]

The formula can be understood by decomposing any vector \(\mathbf v\) into a
component parallel to the axis and a component perpendicular to it:
\[
    \mathbf v_\parallel = (\mathbf a^{\T}\mathbf v)\mathbf a,
    \qquad
    \mathbf v_\perp = \mathbf v-\mathbf v_\parallel.
\]
A rotation about \(\mathbf a\) leaves \(\mathbf v_\parallel\) unchanged and
rotates \(\mathbf v_\perp\) in the plane perpendicular to \(\mathbf a\).  The
matrix above is the compact Cartesian form of that operation.

Only discrete trial angles are tested:
\[
    \theta = \frac{2\pi}{n},
    \qquad
    n=2,3,\dots,n_{\max}.
\]
The integer \(n\) is the candidate rotation order.  For example, \(n=2\)
corresponds to \(180^\circ\), \(n=3\) to \(120^\circ\), \(n=4\) to
\(90^\circ\), and \(n=6\) to \(60^\circ\).  Linear molecules, which may have
continuous rotational symmetry, are handled separately.

\subsection{Mirror Reflections}

Let \(\mathbf n\in\R^3\) be a unit vector normal to a candidate mirror plane.
The mirror plane is the plane through the origin perpendicular to \(\mathbf n\):
\[
    \{\mathbf r\in\R^3 : \mathbf n^{\T}\mathbf r=0\}.
\]
The reflection matrix through this plane is
\[
    \sigma(\mathbf n)
    =
    I - 2\mathbf n\mathbf n^{\T}.
\]

To derive this formula, decompose any vector \(\mathbf r\) into a component
normal to the plane and a component lying in the plane.  Since \(\mathbf n\) is
unit length, the normal component is
\[
    \mathbf r_\perp
    =
    (\mathbf n^{\T}\mathbf r)\mathbf n.
\]
The matrix \(\mathbf n\mathbf n^{\T}\) is the projection matrix onto
\(\mathbf n\), because
\[
    \mathbf n\mathbf n^{\T}\mathbf r
    =
    \mathbf n(\mathbf n^{\T}\mathbf r)
    =
    (\mathbf n^{\T}\mathbf r)\mathbf n
    =
    \mathbf r_\perp.
\]
The in-plane component is
\[
    \mathbf r_\parallel
    =
    \mathbf r-\mathbf r_\perp.
\]
A mirror reflection keeps \(\mathbf r_\parallel\) unchanged and flips
\(\mathbf r_\perp\):
\[
\begin{aligned}
    \mathbf r'
    &=
    \mathbf r_\parallel - \mathbf r_\perp\\
    &=
    (\mathbf r-\mathbf r_\perp)-\mathbf r_\perp\\
    &=
    \mathbf r-2\mathbf r_\perp\\
    &=
    \mathbf r - 2\mathbf n\mathbf n^{\T}\mathbf r\\
    &=
    (I-2\mathbf n\mathbf n^{\T})\mathbf r.
\end{aligned}
\]
Thus \(\sigma(\mathbf n)=I-2\mathbf n\mathbf n^{\T}\).

\subsection{Improper Rotations}

Improper rotations are tested because many molecular point groups contain
orientation-reversing operations.  Examples include mirror planes, inversion,
and roto-reflection axes.  Without improper operations, molecules such as planar
\(D_{2h}\) systems would be classified incompletely.

For a unit axis \(\mathbf a\), an improper rotation can be written as
\[
    S(\mathbf a,\theta)
    =
    R(\mathbf a,\theta)\,\sigma(\mathbf a),
\]
where \(R(\mathbf a,\theta)\) is a proper rotation about \(\mathbf a\), and
\(\sigma(\mathbf a)\) is the reflection through the plane perpendicular to
\(\mathbf a\).  The determinant of a proper rotation is \(+1\), whereas the
determinant of a reflection is \(-1\).  Therefore the determinant of an
improper rotation is \(-1\).

\section{Verification by Atom Permutation}

A trial matrix \(M\) is accepted as a molecular symmetry operation only if it
maps the molecule onto itself.  This is tested by requiring a bijection
\(\pi\) of atom indices such that
\[
    \lVert M\mathbf r_i-\mathbf r_{\pi(i)}\rVert \le \tau
\]
for every atom \(i\), with matching element labels
\[
    e_i = e_{\pi(i)}.
\]
Here:
\begin{itemize}
    \item \(M\) is the trial symmetry matrix;
    \item \(\mathbf r_i\) is the original centered coordinate of atom \(i\);
    \item \(M\mathbf r_i\) is the transformed coordinate;
    \item \(\pi(i)\) is the index of the original atom matched to transformed
    atom \(i\);
    \item \(\tau\) is the distance tolerance.
\end{itemize}

The word ``bijection'' means that the mapping is one-to-one and onto: every
transformed atom must match exactly one original atom, and every original atom
must be used exactly once.  This prevents two transformed atoms from being
assigned to the same original atom, and it prevents any original atom from being
left unmatched.

The candidate-axis construction is therefore allowed to be generous.  It may
generate directions that are not actual symmetry directions.  Such candidates
are harmless, because the final atom-permutation verification rejects any trial
matrix that does not map the full molecule onto itself.

\section{Operation Identity: \((\pi,\det M)\)}

\subsection{The Permutation as an Atom Mapping}

The permutation \(\pi\) is a mapping of atom indices.  If the molecule has
\(N\) atoms, then
\[
    \pi:\{1,\dots,N\}\rightarrow\{1,\dots,N\}.
\]
For a trial symmetry matrix \(M\), \(\pi(i)\) is the index of the atom onto
which atom \(i\) is mapped:
\[
    M\mathbf r_i \approx \mathbf r_{\pi(i)}.
\]
The approximation is controlled by the distance tolerance \(\tau\).  Since a
valid symmetry operation preserves element identity and distance environment,
\(\pi(i)\) must lie in the same equivalence class as \(i\).  The mapping must be
a bijection, so every atom is used exactly once.

\subsection{Why \(\pi\) Alone Is Not Always Enough}

For a rigid nonlinear molecule in a generic three-dimensional arrangement, the
atom permutation often uniquely determines the operation.  However, planar
molecules provide important counterexamples.

Suppose the molecule lies in the \(xy\)-plane, so every atom has \(z_i=0\).  The
mirror through the molecular plane is
\[
    \sigma_{xy}
    =
    \begin{bmatrix}
    1 & 0 & 0\\
    0 & 1 & 0\\
    0 & 0 & -1
    \end{bmatrix}.
\]
For any atom in the plane,
\[
    \sigma_{xy}\mathbf r_i=\mathbf r_i.
\]
Thus the mirror operation induces the identity atom permutation:
\[
    \pi(i)=i.
\]
The identity matrix \(I\) also induces the same atom permutation:
\[
    I\mathbf r_i=\mathbf r_i.
\]
Therefore the permutation \(\pi\) alone cannot distinguish the molecular-plane
mirror from the identity operation.

The determinant resolves this ambiguity:
\[
    \det I = +1,
    \qquad
    \det \sigma_{xy}=-1.
\]
The pair
\[
    (\pi,\det M)
\]
therefore separates proper and improper operations that may induce the same
atom mapping on planar or otherwise geometrically special molecules.

\subsection{Why the Matrix Itself Is Not Used as the Primary Identity}

In exact mathematical geometry, a symmetry operation has a well-defined matrix.
However, the algorithm operates on floating-point molecular coordinates, often
from relaxed structures that are only approximately symmetric.  Candidate axes
are computed from quantities such as
\[
    \mathbf r_i+\mathbf r_j,\qquad
    \mathbf r_i-\mathbf r_j,\qquad
    \mathbf r_i\times\mathbf r_j,
\]
or from eigenvectors of \(T\).  These directions are floating-point vectors.

Two different candidate constructions may describe the same symmetry axis but
with slightly different numerical components.  For example, two candidate axes
that should both represent the \(y\)-axis may be
\[
    \mathbf a_1=(10^{-6},0.999999,2\times10^{-6}),
    \qquad
    \mathbf a_2=(0,1,-10^{-6}).
\]
Even with the same exact trial angle \(\theta=\pi\), the matrices
\[
    R(\mathbf a_1,\pi)
    \qquad\text{and}\qquad
    R(\mathbf a_2,\pi)
\]
will not be bitwise identical.  If matrices were used as operation identities,
one would need a matrix tolerance \(\epsilon\) and compare
\[
    \lVert M_1-M_2\rVert < \epsilon.
\]
Choosing \(\epsilon\) is delicate: too small a tolerance can split one operation
into several duplicates, while too large a tolerance can merge distinct
operations.

By contrast, once a trial operation passes the atom-matching test, its
permutation \(\pi\) is discrete integer data.  Comparing two permutations is an
exact comparison of integer arrays.  The determinant adds the missing parity
information.  Hence \((\pi,\det M)\) is used as the computational identity of an
operation.

\subsection{Deduplication as an Exact Integer Operation}

During candidate testing, the algorithm may rediscover the same symmetry
operation multiple times.  Deduplication means removing repeated operations.

If the operation identity were a floating-point matrix, deduplication would
require a numerical matrix comparison.  With the identity \((\pi,\det M)\),
deduplication becomes exact.  For example, a permutation can be stored as an
integer array:
\[
    \pi=[3,4,1,2],
\]
meaning
\[
    1\mapsto3,\qquad
    2\mapsto4,\qquad
    3\mapsto1,\qquad
    4\mapsto2.
\]
Two operations are duplicates if and only if their integer arrays are identical
and their determinants are identical.  No floating-point tolerance is involved.
This is what is meant by an ``integer operation'' in this context.

\subsection{Group Closure and Permutation Composition}

The molecular symmetry operations form a group.  This means:
\begin{enumerate}[label=(\alph*)]
    \item the identity operation is present;
    \item every operation has an inverse;
    \item the product of any two operations is also an operation;
    \item operation multiplication is associative.
\end{enumerate}
The closure condition is especially useful algorithmically.  If \(g\) and \(h\)
are known symmetry operations, then their product \(gh\) must also be present.

Let
\[
    g=(\pi_g,\det g),
    \qquad
    h=(\pi_h,\det h).
\]
If \(h\) acts first and \(g\) acts second, then the product operation \(gh\)
maps atom \(i\) according to
\[
    \pi_{gh}(i)=\pi_g(\pi_h(i)).
\]
This is called permutation composition.  The determinant of the product is
\[
    \det(gh)=\det(g)\det(h).
\]
Thus group multiplication can be carried out using only integer permutation
composition and multiplication of signs \(\pm1\).

As a concrete example, suppose
\[
    \pi_h=[2,1,3],
    \qquad
    \pi_g=[1,3,2].
\]
Then
\[
    \pi_{gh}(1)=\pi_g(\pi_h(1))=\pi_g(2)=3,
\]
\[
    \pi_{gh}(2)=\pi_g(\pi_h(2))=\pi_g(1)=1,
\]
\[
    \pi_{gh}(3)=\pi_g(\pi_h(3))=\pi_g(3)=2.
\]
Therefore
\[
    \pi_{gh}=[3,1,2].
\]

\subsection{The Group Multiplication Table}

After the operation set is closed, one can store a multiplication table.  If the
operations are indexed as
\[
    g_1,g_2,\dots,g_m,
\]
then the multiplication table records, for each pair \((i,j)\), the index \(k\)
such that
\[
    g_i g_j = g_k.
\]
This table is useful for:
\begin{itemize}
    \item verifying that the detected operation set is closed;
    \item generating missing operations during closure;
    \item classifying the point group;
    \item checking later whether a crystallographic site-symmetry group is a
    subgroup of the molecular symmetry group after orientation.
\end{itemize}

\section{Nonlinear and Linear Molecules}

\subsection{Definitions and Examples}

A molecule is nonlinear if its atoms are not all collinear.  Examples include
water, ammonia, methane, benzene, naphthalene, and PTCDA.

A molecule is linear if all atoms lie on a single straight line.  Examples
include carbon dioxide, acetylene, hydrogen cyanide, and ideal diatomic
molecules such as \(\mathrm{N}_2\) and \(\mathrm{O}_2\).

In terms of the gyration tensor, an ideal linear molecule has molecular spread
only along one direction.  Therefore two eigenvalues of \(T\) are zero, or
nearly zero for numerical geometries.  A nonlinear molecule has at least two
independent directions of atomic spread.

\subsection{Why Linear Molecules Need Separate Handling}

The finite point-group pipeline assumes that the molecule has a finite set of
symmetry operations.  This is true for nonlinear molecular point groups.

An ideal linear molecule has continuous rotational symmetry about the molecular
axis.  For example, an ideal \(\mathrm{CO}_2\) molecule can be rotated by any
angle \(\theta\in[0,2\pi)\) about its molecular axis without changing the
molecule.  This is not a finite set of rotations such as \(C_2\), \(C_3\),
\(C_4\), or \(C_6\).  It is an infinite continuous family.

Linear molecules are therefore classified separately as
\[
    C_{\infty v}
    \qquad\text{or}\qquad
    D_{\infty h},
\]
depending on the presence of inversion or equivalent end-exchange symmetry.
They do not fit the general procedure of enumerating a finite list of candidate
rotation orders and building a finite group multiplication table.

\section{Orthogonal Procrustes Matrix Refinement}

\subsection{The Refinement Problem}

After a trial operation has been accepted, the algorithm knows the atom mapping
\(\pi\).  The raw trial matrix may have been generated from a noisy
floating-point axis.  Matrix refinement replaces that raw matrix with the
orthogonal matrix that best fits the accepted atom mapping.

Let
\[
    \mathbf x_i=\mathbf r_i
    \qquad\text{and}\qquad
    \mathbf y_i=\mathbf r_{\pi(i)}.
\]
Here:
\begin{itemize}
    \item \(\mathbf x_i\) is the original centered coordinate of atom \(i\);
    \item \(\mathbf y_i\) is the centered coordinate of the atom assigned to
    atom \(i\) by the accepted permutation \(\pi\);
    \item \(M\) is the orthogonal matrix to be refined.
\end{itemize}
The orthogonal Procrustes problem is
\[
    M^\star
    =
    \arg\min_{M^{\T}M=I,\ \det M=d}
    \sum_{i=1}^{N}
    \lVert M\mathbf x_i-\mathbf y_i\rVert^2,
\]
where \(d\in\{+1,-1\}\) is the desired determinant.  The determinant constraint
preserves whether the operation is proper (\(d=+1\)) or improper (\(d=-1\)).

\subsection{Cross-Covariance Matrix}

Define the cross-covariance matrix
\[
    H
    =
    \sum_{i=1}^{N}
    \mathbf y_i\mathbf x_i^{\T}.
\]
The matrix \(H\) is \(3\times3\).  It measures how the target coordinates
\(\mathbf y_i\) correlate with the source coordinates \(\mathbf x_i\).  The
normalization factor sometimes used in statistical covariance is irrelevant
here, because multiplying \(H\) by a positive scalar does not change the
maximizing orthogonal matrix.

\subsection{Singular Value Decomposition}

The singular value decomposition of \(H\) is
\[
    H=U\Sigma V^{\T},
\]
where:
\begin{itemize}
    \item \(U\) is a \(3\times3\) orthogonal matrix;
    \item \(V\) is a \(3\times3\) orthogonal matrix;
    \item \(\Sigma=\operatorname{diag}(\sigma_1,\sigma_2,\sigma_3)\);
    \item \(\sigma_1\ge\sigma_2\ge\sigma_3\ge0\) are the singular values.
\end{itemize}
The refined matrix has the form
\[
    M^\star
    =
    UDV^{\T},
\]
where
\[
    D=
    \operatorname{diag}
    \left(1,1,d\,\det(UV^{\T})\right).
\]
This choice guarantees
\[
    \det M^\star=d.
\]

\subsection{Derivation of the Procrustes Solution}

Start from the objective function
\[
    F(M)
    =
    \sum_i
    \lVert M\mathbf x_i-\mathbf y_i\rVert^2.
\]
Expanding one term gives
\[
\begin{aligned}
    \lVert M\mathbf x_i-\mathbf y_i\rVert^2
    &=
    (M\mathbf x_i-\mathbf y_i)^{\T}
    (M\mathbf x_i-\mathbf y_i)\\
    &=
    \mathbf x_i^{\T}M^{\T}M\mathbf x_i
    -
    \mathbf x_i^{\T}M^{\T}\mathbf y_i
    -
    \mathbf y_i^{\T}M\mathbf x_i
    +
    \mathbf y_i^{\T}\mathbf y_i.
\end{aligned}
\]
The two cross terms are equal because they are scalar transposes of one another:
\[
    \mathbf x_i^{\T}M^{\T}\mathbf y_i
    =
    \left(\mathbf x_i^{\T}M^{\T}\mathbf y_i\right)^{\T}
    =
    \mathbf y_i^{\T}M\mathbf x_i.
\]
Since \(M^{\T}M=I\),
\[
    \mathbf x_i^{\T}M^{\T}M\mathbf x_i
    =
    \mathbf x_i^{\T}\mathbf x_i.
\]
Therefore
\[
    F(M)
    =
    C
    -
    2\sum_i\mathbf y_i^{\T}M\mathbf x_i,
\]
where
\[
    C=
    \sum_i
    \left(
    \lVert\mathbf x_i\rVert^2+
    \lVert\mathbf y_i\rVert^2
    \right)
\]
does not depend on \(M\).  Minimizing \(F(M)\) is therefore equivalent to
maximizing
\[
    \sum_i\mathbf y_i^{\T}M\mathbf x_i.
\]
Using
\[
    H=\sum_i\mathbf y_i\mathbf x_i^{\T},
\]
we can write
\[
    \sum_i\mathbf y_i^{\T}M\mathbf x_i
    =
    \operatorname{tr}(MH^{\T}).
\]
With \(H=U\Sigma V^{\T}\), we have \(H^{\T}=V\Sigma U^{\T}\).  Hence
\[
\begin{aligned}
    \operatorname{tr}(MH^{\T})
    &=
    \operatorname{tr}(MV\Sigma U^{\T})\\
    &=
    \operatorname{tr}(U^{\T}MV\Sigma),
\end{aligned}
\]
where the second equality uses cyclic invariance of the trace.

Define
\[
    A=U^{\T}MV.
\]
Since \(U\), \(M\), and \(V\) are orthogonal, \(A\) is also orthogonal.  The
problem becomes
\[
    \max_{A^{\T}A=I,\ \det A=\delta}
    \operatorname{tr}(A\Sigma),
\]
where
\[
    \delta=d\,\det(UV^{\T}).
\]
The determinant relation follows from
\[
    \det A
    =
    \det(U^{\T})\det(M)\det(V)
    =
    d\,\det(U)\det(V)
    =
    d\,\det(UV^{\T}).
\]

Because \(\Sigma\) is diagonal,
\[
    \operatorname{tr}(A\Sigma)
    =
    \sigma_1 A_{11}+\sigma_2 A_{22}+\sigma_3 A_{33}.
\]
For an orthogonal matrix, the entries cannot exceed one in magnitude.  The
maximum is obtained by aligning the largest singular values with positive
diagonal entries.  If \(\delta=+1\), the best choice is
\[
    A^\star=I.
\]
If \(\delta=-1\), one sign must be negative; because
\(\sigma_3\) is the smallest singular value, the least costly negative sign is
placed on the third axis:
\[
    A^\star=\operatorname{diag}(1,1,-1).
\]
Both cases are written compactly as
\[
    A^\star
    =
    D
    =
    \operatorname{diag}(1,1,\delta).
\]
Since \(A=U^{\T}MV\), this gives
\[
    U^{\T}M^\star V=D,
\]
and therefore
\[
    M^\star=UDV^{\T}.
\]
Substituting the definition of \(\delta\) yields
\[
    \boxed{
    M^\star
    =
    U
    \operatorname{diag}
    \left(1,1,d\,\det(UV^{\T})\right)
    V^{\T}
    }.
\]
This is the least-squares best orthogonal matrix with the prescribed
determinant for the accepted atom mapping.

\subsection{Why Refinement Is Needed}

The raw trial matrix is generated from candidate axes and trial angles.  The
axis may be slightly noisy because it is computed from floating-point molecular
coordinates.  Once the atom mapping \(\pi\) is known, the Procrustes step uses
all atoms simultaneously to compute the best orthogonal matrix consistent with
that mapping.  The result is a cleaner operation matrix, while the operation
identity remains the exact discrete pair \((\pi,\det M)\).

\section{Recording Symmetry Elements}

\subsection{Operation Identity Versus Geometric Metadata}

The pair \((\pi,\det M)\) is useful for exact computation, but it is not a full
human-readable geometric description.  Therefore each accepted operation should
store both:
\begin{enumerate}[label=(\alph*)]
    \item the computational identity \((\pi,\det M)\);
    \item geometric metadata extracted from the refined matrix \(M^\star\).
\end{enumerate}
The stored data include:
\begin{itemize}
    \item the refined Cartesian matrix \(M^\star\);
    \item the atom permutation \(\pi\);
    \item the determinant \(\det M^\star\);
    \item the maximum atom-matching deviation;
    \item a classified operation type, such as \(E\), \(i\), \(C_n\),
    \(\sigma\), or \(S_n\).
\end{itemize}

\subsection{Proper Rotation Elements}

For a proper rotation,
\[
    \det M^\star=+1.
\]
The rotation angle \(\theta\) can be obtained from the trace:
\[
    \operatorname{tr}(M^\star)=1+2\cos\theta,
\]
so
\[
    \cos\theta=
    \frac{\operatorname{tr}(M^\star)-1}{2}.
\]
The rotation axis \(\mathbf a\) is the eigenvector with eigenvalue \(1\):
\[
    M^\star\mathbf a=\mathbf a.
\]
If
\[
    \theta\approx\frac{2\pi}{n},
\]
the element is recorded as a \(C_n\) axis along \(\mathbf a\).

\subsection{Mirror, Inversion, and Improper Elements}

For a pure mirror operation, the eigenvalues are
\[
    1,\quad 1,\quad -1.
\]
The mirror normal \(\mathbf n\) is the eigenvector with eigenvalue \(-1\):
\[
    M^\star\mathbf n=-\mathbf n.
\]
The mirror plane is
\[
    \{\mathbf r:\mathbf n^{\T}\mathbf r=0\}.
\]

The inversion operation is identified by
\[
    M^\star=-I.
\]
It is recorded as inversion through the molecular centroid, which is the origin
of the centered coordinate system.

An improper rotation is recorded as an \(S_n\) element.  It has determinant
\(-1\) and can be interpreted as a proper rotation about an axis followed by
reflection through the plane perpendicular to that axis.  The axis and order are
obtained from the refined matrix and, when useful, from the trial operation that
produced the accepted atom mapping.

\section{Classifying the Point Group}

After all valid operations have been found, deduplicated, and closed under
group multiplication, the point group symbol is determined from the inventory of
symmetry elements.

\subsection{Linear Molecules}

Linear molecules are classified first, because they have continuous rotational
symmetry and do not fit the finite-operation classification:
\[
    C_{\infty v}
    \qquad\text{or}\qquad
    D_{\infty h}.
\]

\subsection{Cubic Families}

For nonlinear molecules, the algorithm first checks for multiple high-order
proper rotation axes.  If there are two or more distinct axes of order at least
three, the molecule belongs to a cubic family:
\[
    T,\quad T_d,\quad T_h,\quad O,\quad O_h,\quad I,\quad I_h.
\]
The distinction is made from the set of rotation orders and the presence or
absence of inversion and mirror-type improper elements.

\subsection{\(C_n\) and \(D_n\) Families}

If the molecule is not cubic, the highest-order proper rotation axis is chosen
as the principal axis.  Let its order be \(n\).

If there are \(n\) two-fold axes perpendicular to the principal axis, the group
belongs to the \(D_n\) family:
\[
    D_n,\qquad D_{nh},\qquad D_{nd}.
\]
The suffix is determined by improper elements:
\begin{itemize}
    \item a horizontal mirror plane gives \(D_{nh}\);
    \item suitable diagonal mirror planes give \(D_{nd}\);
    \item neither gives \(D_n\).
\end{itemize}

If there are not \(n\) perpendicular \(C_2\) axes, the group belongs to the
\(C_n\) family:
\[
    C_n,\qquad C_{nv},\qquad C_{nh},\qquad S_{2n}.
\]
The suffix is again determined by improper elements:
\begin{itemize}
    \item vertical mirror planes give \(C_{nv}\);
    \item a horizontal mirror plane gives \(C_{nh}\);
    \item an improper axis without the corresponding proper-axis structure gives
    an \(S_{2n}\) group;
    \item no relevant improper element gives \(C_n\).
\end{itemize}

\subsection{No Nontrivial Proper Rotation}

If no nontrivial proper rotation axis is present, the remaining possibilities
are:
\[
    C_s,\qquad C_i,\qquad C_1.
\]
Specifically:
\begin{itemize}
    \item mirror only gives \(C_s\);
    \item inversion only gives \(C_i\);
    \item identity only gives \(C_1\).
\end{itemize}

For the planar molecules expected in the example systems, such as naphthalene
and PTCDA, the expected classification is \(D_{2h}\): they have three mutually
perpendicular \(C_2\) axes together with inversion and mirror symmetry.

\end{document}
